\begin{center}
  {\small{\em À medida que você simplifica sua vida, as leis do universo se tornam mais simples.} -- Henry David Thoreau}
\end{center}

Há algum tempo atrás, ouvi a história de alguns garotos que entraram à noite num supermercado e, só por diversão, trocaram todas as etiquetas de preço. As etiquetas de preço das bananas foram colocadas nas TVs de última geração, e vice-versa. Uma geladeira passou a custar menos de um real, enquanto um litro de leite passava dos quinhentos reais. O resultado, no dia seguinte, foi obviamente o caos.

O mesmo aconteceu com nosso mundo ocidental moderno. Alguém, há algum tempo atrás, trocou as etiquetas de preço das coisas deste mundo. Coisas sem valor nenhum passaram a se tornar valiosas aos olhos das pessoas, enquanto o que tem valor real passou a ser desprezado.

Esse é o mundo em que nascemos, eu e você. Um mundo complicado. Um mundo tão complicado que não sabemos nem por onde começar para compreendê-lo. Na verdade, não precisamos tentar entendê-lo. Dentro de cada casa existe um aparelho chamado televisão que nos ensina como o mundo funciona, no que devemos acreditar, o que devemos comprar, e quais devem ser os nossos valores. E nós obedecemos.

O problema é que estes valores que somos ensinados estão todos invertidos. E a não ser que de alguma forma sejamos sacudidos para fora desta visão de mundo e passemos a questionar aquilo que vemos e aprendemos, é possível que passemos nossas vidas inteiras acreditando em valores errados.

Alguns dos valores que nossa cultura nos ensina são:

\textbf{O valor de um homem se mede pela quantidade de bens que ele adquiriu.} Este tipo de valor nunca é medido em isolamento, mas sempre em comparação. Um homem nunca é rico ou pobre sozinho, mas sempre comparado aos seus vizinho, amigos, parentes ou colegas. E é por isso que a noção de riqueza e pobreza varia grandemente de um país para o outro.

Essa idéia de valor pessoal costuma ultrapassar qualquer outra -- assim, não é incomum vermos pessoas gastando tudo o que tem (e o que não tem) na tentativa de parecerem possuir mais condições do que realmente tem. Um corretor imobiliário em Hollywood, um dos metros quadrados mais caros do planeta, certa vez disse que muitas pessoas alugam mansões sem ter dinheiro para mobiliá-las. Assim, vivem dentro delas sem móveis nem condições, apenas com o intuito de serem valorizadas perante seus vizinhos.

E assim, como disse Thoreau, {\em``a civilização tem melhorado suas casas, mas não tem melhorado igualmente os homens que habitam nelas. Tem criado palácios, mas não tem sido assim tão fácil criar nobres e reis.''}

A idéia do valor pessoal que depende da quantidade de bens é reforçada diariamente pelos programas e propagandas na televisão. Assim, um homem sente que seu \emph{status} está em jogo não apenas perante seus pares, mas também perante sua família. Sua própria visão de valor pessoal passa a depender de sua capacidade de suprir os desejos de consumo de sua família, seja quantas forem as horas necessárias de trabalho.

Provavelmente você conheça pessoas assim. E antes de julgá-las, avalie sua própria vida. Você se sente rico ou pobre? Quanto tempo diário você tem gasto na aquisição de recursos? Com que propósito?


\textbf{A televisão provê um descanso merecido após um dia de trabalho.} De preferência no sofá, sem camisa e com uma cerveja. Não é essa imagem que a própria televisão passa?

\textbf{Compre agora. Pague depois.}

\textbf{O sistema provê tudo que necessitamos, sejam necessidades ou desejos.}

\textbf{A verdadeira felicidade está na obtenção da gratificação imediata.}

\textbf{Vamos viver para sempre.}


